\section{The Memory-Augmented Agent}
\label{sec:agent}

\begin{figure*}[t]
\centering
\begin{tikzpicture}[
  scale=0.70, transform shape,
  font=\footnotesize,
  >={Latex[length=2mm]},
  term/.style={draw, rounded corners, align=center, minimum width=23mm, minimum height=11mm, fill=black!4},
  core/.style={draw, very thick, rounded corners, align=center, minimum width=26mm, minimum height=11mm, fill=black!8},
  mem/.style={draw, cylinder, shape border rotate=90, aspect=0.16, align=center,
              minimum width=27mm, minimum height=11mm, fill=black!3, inner sep=1pt},
  llm/.style={draw, dashed, rounded corners, align=center, minimum width=25mm, minimum height=12mm, fill=black!2},
  note/.style={font=\scriptsize\itshape, align=center, text=black!65},
  el/.style={font=\scriptsize, align=center},
]

% --- nodes ---
\node[term] (term) at (0,0) {\textbf{Terminal}\\ \textbf{multiplexer}\\[1pt] \scriptsize editor + shell panes};
\node[core] (core) at (5.4,0) {\textbf{Agent}\\ \textbf{core}};
\node[mem]  (m1) at (3.35,-2.35) {\textbf{M1: your commands}\\[1pt] \scriptsize \code{\char`~/.mars/cmd\_memory}\\ \scriptsize + shell history};
\node[mem]  (m2) at (7.45,-2.35) {\textbf{M2: tool self-knowledge}\\[1pt] \scriptsize docs, action registry,\\ \scriptsize tuning knobs};
\node[llm]  (llm) at (12.7,0) {\textbf{LLM}\\ model\\ cascade};
\node[note] (log) at (5.4,1.55) {self-instrumenting log\\ $=$ benchmark corpus};

% --- your-machine boundary (behind) ---
\begin{scope}[on background layer]
  \node[draw, dashed, thick, rounded corners=6pt, fill=blue!3, inner sep=10pt,
        fit=(term)(core)(m1)(m2)(log)] (mach) {};
\end{scope}
\node[anchor=north west, font=\scriptsize\bfseries, text=blue!45] at (mach.north west) {~your machine};

% --- terminal <-> core ---
\draw[->] ([yshift=2.4mm]term.east) -- ([yshift=2.4mm]core.west)
      node[midway, above, el] {NL request / question};
\draw[<-] ([yshift=-2.4mm]term.east) -- ([yshift=-2.4mm]core.west)
      node[midway, below, el] {directive:\\ \code{RUN}/\code{TYPE}/\code{OPEN}\\ (registry-validated)};
\draw[->, dashed] (term.north) to[bend left=22] node[midway, above, note] {sees the panes (sight)} (core.north west);

% --- core <-> memory (BM25) ---
\draw[<->] (core.south) -- (m1.north);
\draw[<->] (core.south) -- (m2.north);
\node[el] at (5.4,-1.35) {BM25 retrieve\\ $\rightarrow$ few-shot into prompt};

% --- core <-> LLM (the only thing that leaves) ---
\draw[->, very thick] ([yshift=2.4mm]core.east) -- ([yshift=2.4mm]llm.west)
      node[midway, above, el] {prompt};
\draw[<-, very thick] ([yshift=-2.4mm]core.east) -- ([yshift=-2.4mm]llm.west)
      node[midway, below, el] {completion};
\node[note, text=red!55] at (10.05,-1.05) {the only call that\\ leaves your machine};

% --- broker annotation ---
\node[note] (brk) at (12.7,-2.2) {\textbf{key-never-leaves-home}\\ on a remote box the call is\\ proxied home; the key stays local};
\draw[->, dotted, black!55] (brk.north) -- (llm.south);

\end{tikzpicture}
\caption{MARS retrieves with lexical BM25 from two local stores (M1, the user's past
commands; M2, the tool's own docs and registry) and injects the results as few-shot
exemplars. Only the model call leaves the machine; on a remote box even the key stays
home, and every call is logged.}
\label{fig:arch}
\end{figure*}


\subsection{Design considerations}
\label{sec:design}

\emph{The agent should be another input device, not a second system} --- not a
subsystem bolted beside the editor with its own command language, state, and failure
modes. Everything MARS can do is a variant of a single \code{Action} enum, and three
interfaces resolve a user's intent to it: a direct keychord, a fuzzy
command bar, and the agent (natural language). All three terminate in a
single dispatch point, so a capability added once is chord-bindable, searchable, and
agent-invokable. Figure~\ref{fig:arch} sketches the system and its retrieval loop.

This choice trades expressive freedom for grounding: the agent can only name actions
that exist in the registry and retrieves from flat local files, so the surface a user
must trust is small and easy to audit, but MARS cannot invent a capability at runtime,
and lexical retrieval will miss paraphrases an embedding model would catch
(\S\ref{sec:eval}). In the terminal domain we believe this trade is strongly favorable.

\subsection{The agent as a validated input device}
\label{sec:directives}

\emph{An agent that acts on the user's workspace must be no less safe than a keypress,
never a bypass around the editor's gates.} Following the mode-error tradition in interface
design \citep{norman1988design}, we treat the agent as one input device among several,
behind the same confirmation gates as a keystroke. The implementation is a trailing-line
\emph{directive} protocol: the model receives a live catalog of every action,
description, and keybinding, and ends its reply with at most one directive, each gated
by confirmation --- \code{RUN: <Action>} fires a registry action, \code{TYPE: <cmd>}
types a command into a shell pane, \code{OPEN: path:line} jumps to a file location, and
\code{NEED: <context>} pulls more context. We chose plain text over provider-native
tool-calling APIs because it works identically across OpenAI-compatible endpoints, a
native Anthropic path, and local open models (which also keep prompts on the machine;
see Ethics Statement), and renders as human-readable text the user confirms before
anything fires; the cost is the stricter typing a schema provides.

The grounding comes from the registry check: a \code{RUN:} whose name does not resolve
against the live \code{Action} set is rejected before dispatch, so the model cannot
execute a hallucinated capability, and a destructive action passes the same
confirmation gate a human would. The \code{NEED:} directive closes a different gap:
rather than receiving the full screen with every request, the model can ask for more
context (a pane's scrollback, another tab); MARS supplies it and re-asks exactly once,
so no unbounded retrieval loop can form.

\subsection{Natural-language shell translation}
\label{sec:translate}

The bread-and-butter agentic request at a terminal is to turn English into the right
command. \emph{Translation should be assistive and reversible, never an automatic
execution.} The implementation lives on the command bar: the user presses
\code{C-Space} (we write keybindings in Emacs notation: \code{C-} is Ctrl, \code{M-}
is Meta) and types plain English; MARS sends the request with on-screen context
(working directory, recent output) to the model cascade (\S\ref{sec:cascade}) and
renders the proposed command at the cursor to run, edit, or discard. The same translator serves as the command bar's fallback
(\S\ref{sec:system}). Translation by itself is a commodity
capability \citep{lin2018nl2bash,lin2017tellina}; what personalizes it is the memory we
describe next.

\subsection{Memory 1: the user's own commands}
\label{sec:m1}

\emph{The user's own accepted commands are the highest-value context available, and they
are local to the terminal by construction.} Asked to ``launch the training run,'' a
generic model produces a plausible
\code{python train.py} when this project in fact uses
\code{sbatch -{}-gres=gpu:8 scripts/train.slurm}. The implementation is a corrective
loop. When the user accepts (or edits and runs) a translated command, MARS appends the
\code{(request, command)} pair to a local JSON-lines file; it also reads recent shell
history from \code{\$HISTFILE}. On the next request, MARS ranks these entries against the
query with BM25 \citep{robertson2009bm25}, with standard parameters ($k_1{=}1.5$,
$b{=}0.75$) and no index to build, and injects the top three pairs as few-shot exemplars,
together with up to three ranked lines of recent shell history.

The loop is corrective in the precise sense that a command the model got wrong once
becomes one it gets right thereafter, because the user's correction is what enters the
store; a per-invocation assistant never observes the correction. The retriever is
deliberately unremarkable; the finding is that in this domain unremarkable retrieval
is decisive (\S\ref{sec:eval}).

\subsection{Memory 2: the tool's own documentation}
\label{sec:m2}

\emph{An embedded agent should answer questions about its tool from ground truth rather
than a guess, and a registry generated from the code is ground truth that cannot name a
capability that does not exist.} MARS's action registry,
keybindings, and tuning knobs each carry a human-written description, so documentation
is generated from the code itself. The same BM25 retriever runs over this second
corpus (documentation chunked into passages, plus the live registry and described
knobs); on a \code{NEED: docs} pull, MARS injects the top passages, so ``how do I
split the screen?'' is answered with the real keybinding and ``make the accent
orange'' resolves to the actual knob and a confirmable edit. When
retrieval returns nothing, the agent reports that it cannot find such a capability
rather than inventing one.

\subsection{Simplicity as a control}
\label{sec:whysimple}

The instinct when building a memory system is to reach for embeddings and a vector
database; we keep the retriever deliberately trivial as a control, so any gap it
closes is attributable to the information, not the machinery. The terminal also favors lexical retrieval: command
vocabulary is small, repetitive, and lexically stable, so a query and its useful
memory tend to share surface tokens (the regime where BM25 is strong and dense
retrieval's paraphrase advantage is marginal); the corpora are small; and a single
relevant exemplar supplies the exact string the user wants reproduced. We therefore
treat embeddings as a cost to be deferred until the domain broadens
(\S\ref{sec:future}). The empirical support is \S\ref{sec:eval}, in which BM25
bridges paraphrases whose mean word-Jaccard similarity to the query is only 0.36, with
nothing more than a prompt, a local file, and a standard ranking function.

\subsection{The model cascade}
\label{sec:cascade}

\emph{Model choice should be right-sized per task.} Most LLM traffic at a terminal is
cheap and can be served by small or free models, but a single fixed model either wastes
capability on easy tasks or fails hard ones, and hits rate limits where volume is high.
We distinguish two
orthogonal reasons to change models. Rotating \emph{for limits} moves a task
horizontally on a rate-limit error; because each provider meters usage on its own
counter, rotating across $N$ families multiplies the effective free ceiling by roughly
$N$. Escalating \emph{for quality} moves a task vertically to a stronger tier, and only
on a detected failure.
Detection is free for one class of failures: because the agent's output is structured, a
\code{RUN:} that fails the registry check is an unambiguous signal that a small model
erred, and it triggers one retry at a higher tier. The benefit is headroom against
the provider quota walls we later hit during evaluation (\S\ref{sec:setup}).
